\uuid{VFHz}
\titre{Congruences}
\auteur{bodin ; croizat ; sacre}
\organisation{exo7}
\licence{CC BY-NC-SA 4.0}
\datecreate{2022-03-08}
\niveau{L1}
\module{Algèbre}
\chapitre{Arithmétique dans Z}
\sousChapitre{Autre}
\difficulte{3}
\theme{Congruences}
\exo7qcmid{hilisit.arithmetique.049}

\contenu{
\begin{question}
\qclassification{103.99}
    Quelles sont les affirmations vraies ?
    \begin{answers}
        \bad{$2^8 \equiv 2 \;[8]$}
        \bad{$3^{12} \equiv 3 \;[13]$}
        \bad{$18^7 \equiv 1 \; [19]$}
        \good{$4^{16} \equiv 1 \; [17]$}
    \end{answers}
    \begin{explanations}
    $8$ n'est pas un nombre premier, le petit théorème de Fermat ne s'applique pas. En fait $2^8 \equiv 0 \;[8]$ car $2^3 = 8 \equiv 0 \;[8]$. \\
    Petit théorème de Fermat, avec $p=13$,  $3^{12} \equiv 1 \;[13]$. \\
    $18^7 \equiv (-1)^7 \equiv -1  \equiv 18 \; [19]$ (le calcul n'a rien à voir avec le petit théorème de Fermat). \\
    Petit théorème de Fermat, avec $p=17$,  $4^{16} \equiv 1 \; [17]$.
    \end{explanations}
\end{question}
}
